% NON-CANONICAL DRAFT / DO NOT CITE / Lean-verified core: EpiFormal/KappaSharp.lean
% Internal source map (not for submission): Prop 2 = THM-0006; Cor 3 = THM-0003(B);
% Thm 4 = THM-0005; prop:input (Section sec:input) = R5 paper proof of the former "input"
% (finite checks: logs/verification/R5/input_prop_chk.py/.out).
% Citation data only from logs/evidence/Q-R2-05_* and Q-R3-01_*.
\documentclass[11pt]{amsart}
\usepackage[utf8]{inputenc}
\usepackage[T1]{fontenc}
\usepackage{lmodern}
\usepackage{amsmath,amssymb}
\usepackage{geometry}
\geometry{margin=1.1in}
\emergencystretch=3em

\newtheorem{theorem}{Theorem}
\newtheorem{proposition}[theorem]{Proposition}
\newtheorem{corollary}[theorem]{Corollary}
\newtheorem{lemma}[theorem]{Lemma}
\theoremstyle{remark}
\newtheorem{remark}[theorem]{Remark}

\newcommand{\Z}{\mathbb{Z}}
\newcommand{\R}{\mathbb{R}}
\newcommand{\Q}{\mathbb{Q}}
\newcommand{\dist}[1]{\lVert #1\rVert}
\newcommand{\LL}{\mathcal{L}}

\begin{document}

\title[A sharp constant along denominators associated with $e$]{A sharp constant for inhomogeneous approximation along a sequence of denominators associated with $e$}

% TODO before submission: author, affiliation, e-mail, acknowledgements, MSC codes, keywords.
\author{[Author name]}
% TODO before submission
\date{Draft of \today. Non-canonical; not for citation.}

\begin{abstract}
Let $Q_j$ be defined by $Q_0=Q_1=1$ and $Q_{j+1}=(4j+2)Q_j+Q_{j-1}$, and let
$\LL=\Z+\Z\,\frac{e+1}{2}$. For every real $z\notin\LL$ we show
$\limsup_j j\,\dist{Q_j z}\ge 1/4$, while for $z\in\LL$ one has
$j\,\dist{Q_jz}\to0$. We exhibit an explicit $z^*\notin\LL$ with
$j\,\dist{Q_jz^*}\to 1/4$ along even $j$ and $\to0$ along odd $j$, so for this sequence the constant
$1/4$ is optimal and attained. The qualitative dichotomy already follows from a theorem of
Kraaikamp and Liardet; the sharp constant is, to our knowledge, not in the literature we consulted
(a negative search only). The lower bound is deduced from a general dichotomy for
integer sequences satisfying a three-term recurrence with partial quotients tending to
infinity. We also give a self-contained proof that $P_j/Q_j\to e$ and that the $Q_j$ are the
regular continued-fraction denominators of $(e-1)/2=[0;1,6,10,14,\dots]$, so that the results
apply to the setting of Kraaikamp and Liardet. The lower-bound and the optimality statements are machine-checked in Lean~4 (full CI build, run 45759959535);
the even/odd sub-limits are not. We state explicitly what is not claimed.
\end{abstract}

\maketitle

\section{Introduction}

For $x\in\R$ write $\dist{x}$ for the distance from $x$ to the nearest integer. Let $\alpha$
be irrational with continued-fraction denominators $q_n$, and set
$\theta_n=q_n\dist{q_n\alpha}$. Kraaikamp and Liardet \cite{KL} proved a rigidity
statement for the approximation of an arbitrary real $\beta$ along the sequence $(q_n)$:
$\beta\in\alpha\Z+\Z$ if and only if $\dist{q_n\beta}\le\theta_n/4$ for all large $n$
(their Theorem~1), and $\beta\in\alpha\Z+\Z$ if and only if $\dist{q_n\beta}=o(\theta_n)$
(their Theorem~1$'$). Earlier, Larcher \cite{Larcher} treated the bounded-type case
($\dist{q_n\beta}\to0$ iff $\beta\in\Z\alpha+\Z$) and gave a counterexample when the partial
quotients are unbounded. The result of \cite{KL} is quoted in a related form in \cite{BDS}
(Lemma~4.1 there, attributed to \cite{KL}). KL also attribute a related statement for bounded type to
Stewart \cite{Stewart}; we know that statement only through the remark in \cite{KL}, and
\emph{the full text of \cite{Stewart} was not consulted}.
As far as the full texts we read go, \cite{KL} contains no sharpness example and no claim
that the constant $1/4$ is optimal.

\medskip
\noindent\textbf{The sequence.} Define integers by
\begin{gather*}
Q_0=Q_1=1,\qquad Q_{j+1}=(4j+2)Q_j+Q_{j-1}\quad(j\ge1),\\
P_0=1,\ P_1=3,\qquad P_{j+1}=(4j+2)P_j+P_{j-1}\quad(j\ge1).
\end{gather*}
Here the recurrence \emph{is the definition}. The proof of Corollary~\ref{cor:lower} uses that
$P_j/Q_j\to e$; this is proved in Proposition~\ref{prop:input}(i) (Section~\ref{sec:input}) via a
Hermite-type integral. Proposition~\ref{prop:input}(ii) shows moreover that the $Q_j$ are exactly the
regular continued-fraction denominators of $\alpha=(e-1)/2=[0;1,6,10,14,\dots]$. We give a
self-contained proof and do not rely on a reference for this expansion.
Consequently $\LL=\Z+\Z\frac{e+1}{2}=\alpha\Z+\Z$, and our results are statements about the
setting of \cite{KL}. Apart from $P_j/Q_j\to e$, none of the proofs in
Sections~\ref{sec:dichotomy}--\ref{sec:sharp} use the identification.

\begin{theorem}[Main theorem]\label{thm:main}
Let $\LL=\Z+\Z\frac{e+1}{2}$ and let $Q_j$ be as above.
\begin{enumerate}
\item[(a)] If $z\in\LL$ then $j\,\dist{Q_jz}\to0$.
\item[(b)] If $z\notin\LL$ then $\limsup_{j\to\infty} j\,\dist{Q_jz}\ge 1/4$.
\item[(c)] There is an explicit $z^*\notin\LL$ (Section~\ref{sec:sharp}) with
$j\,\dist{Q_jz^*}\to 1/4$ along even $j$ and $j\,\dist{Q_jz^*}\to0$ along odd $j$.
\end{enumerate}
Consequently $\kappa:=\inf_{z\notin\LL}\limsup_j j\,\dist{Q_jz}=1/4$, and the infimum is attained.
\end{theorem}

\begin{remark}[Comparison with \cite{KL}]
By Proposition~\ref{prop:input}(ii), $\theta_j=Q_j\dist{Q_j\alpha}$, and the standard identity
$\theta_n=1/(x_{n+1}+q_{n-1}/q_n)$ (with $x_{n+1}$ the complete quotient, $c_{n+1}<x_{n+1}<c_{n+1}+1$)
gives $\theta_j(4j+2)\to1$. So the
bound $\theta_j/4$ in KL's Theorem~1 corresponds, in the normalisation $j\,\dist{Q_jz}$, to a
threshold tending to $1/16$. Theorem~\ref{thm:main}(b) raises this to $1/4$, and (c)
shows that $1/4$ cannot be improved along this sequence. Qualitatively, (a) and (b) are
consequences of \cite{KL} (Theorem~1$'$);
the new content, if any, is the numerical value of the optimal constant and its attainment. Numerically, $\theta_j(4j+2)\to1$ and $j\theta_j/4$ increases to $1/16$ ($0.0594558$ at $j=10$,
$0.0623050$ at $j=160$; mpmath, 900 digits, script \texttt{logs/verification/R4/note\_chk.py});
this is a finite computation, not a proof.
\end{remark}

\medskip
\noindent\textbf{What is not claimed.} See Section~\ref{sec:scope}: nothing here concerns $\pi$,
$e+\pi$ or any specific constant other than $e$ via the lattice $\LL$; sharpness is shown
for this single sequence only; novelty rests on a negative literature search.

\section{A general dichotomy}\label{sec:dichotomy}

Let $\|\cdot\|$ be as above. We use the following hypotheses.

\medskip
\noindent\textbf{(H)} $a_j\in\Z_{>0}$ for $j\ge1$ with $a_j\to\infty$; integers $Q_0\ge0$, $Q_1\ge1$,
$P_0,P_1$; $D:=P_0Q_1-P_1Q_0\ne0$; and for $j\ge1$,
\[
Q_{j+1}=a_jQ_j+Q_{j-1},\qquad P_{j+1}=a_jP_j+P_{j-1}.
\]

\begin{proposition}\label{prop:dich}
Assume (H). Then $\theta:=\lim_j P_j/Q_j$ exists and is irrational, and
$\delta_j:=Q_j\theta-P_j$ satisfies $|\delta_j|\le|D|/Q_{j+1}$. Put
$\LL_D:=(\Z\delta_0+\Z\delta_1)/D$; it contains $\Z+\Z\theta$ (with index $|D|$ in the sense that
$\LL_D$ is the image of $M^{-1}\Z^2\supseteq\Z^2$, $|\det M|=|D|$, see below). For every real $z$:
\begin{enumerate}
\item[(A)] if $z\in\LL_D$ then $a_j\dist{Q_jz}\to0$;
\item[(B)] if $z\notin\LL_D$ then $\limsup_j a_j\dist{Q_jz}\ge1$;
\item[(C)] if $w_j>0$ and $w_j/a_j\to\kappa\in(0,\infty)$, then (A) and (B) hold with $a_j$
replaced by $w_j$ and the threshold $1$ replaced by $\kappa$.
\end{enumerate}
\end{proposition}

\begin{proof}
\emph{Preliminaries.}
(a) Put $C_j:=P_jQ_{j+1}-P_{j+1}Q_j$. By the recurrences,
$C_j=P_j(a_jQ_j+Q_{j-1})-(a_jP_j+P_{j-1})Q_j=P_jQ_{j-1}-P_{j-1}Q_j=-C_{j-1}$,
so $C_j=(-1)^jD$.

(b) Since $a_j\ge1$, $Q_1\ge1$ and $Q_0\ge0$, we have $Q_j\ge1$ for $j\ge1$ and
$Q_{j+1}\ge a_jQ_j$; as $a_j\to\infty$, $Q_j\to\infty$. Also $Q_{k+2}>Q_k$ for $k\ge0$,
so $Q_kQ_{k+1}$ is strictly increasing for $k\ge1$.

(c) From (a), $\frac{P_{k+1}}{Q_{k+1}}-\frac{P_k}{Q_k}=-\frac{(-1)^kD}{Q_kQ_{k+1}}$. For $k\ge1$ these
terms alternate in sign and decrease in absolute value to $0$ by (b). Hence $\theta$ exists, and
$|\theta-P_j/Q_j|\le|D|/(Q_jQ_{j+1})$ for $j\ge1$, so $|\delta_j|\le|D|/Q_{j+1}$.
For $j=0$: if $Q_0=0$ then $|\delta_0|=|P_0|=|D|/Q_1$; if $Q_0\ge1$ the same alternating-series
argument applies from $k=0$ because $Q_2>Q_0$.

(d) \emph{$\theta$ is irrational.} If $\theta=p/q$, then $q\delta_j=pQ_j-qP_j\in\Z$ and $q\delta_j\to0$, so
$\delta_j=0$ for $j\ge N$. The recurrence $\delta_{j-1}=\delta_{j+1}-a_j\delta_j$ (note $\delta$
satisfies the same recurrence as $Q,P$) run backwards gives $\delta_0=\delta_1=0$, hence
$P_i=\theta Q_i$ for $i=0,1$, so $D=P_0Q_1-P_1Q_0=\theta(Q_0Q_1-Q_1Q_0)=0$, a contradiction.

(e) A direct computation gives $Q_j\delta_{j+1}-Q_{j+1}\delta_j=C_j=(-1)^jD\ne0$. Hence any real
sequence satisfying the recurrence (for $j\ge j_0$) is a unique real combination
$\alpha Q_j+\beta\delta_j$ of the two solutions $(Q_j)$ and $(\delta_j)$.

\emph{Lemma 1.} Let $\gamma,\varepsilon\in\R$. If $\gamma Q_j+\varepsilon P_j\in\Z$ for two
consecutive indices $j_0,j_0+1$, then it is an integer for all $j\ge0$. Indeed, the sequence
$u_j=\gamma Q_j+\varepsilon P_j$ satisfies $u_{j+1}=a_ju_j+u_{j-1}$ with integer $a_j$, which
propagates integrality forwards and (via $u_{j-1}=u_{j+1}-a_ju_j$) backwards.
Consequently
$\Lambda:=\{(\gamma,\varepsilon)\in\R^2:\gamma Q_j+\varepsilon P_j\in\Z\ \forall j\}=M^{-1}\Z^2$
with $M=\begin{pmatrix}Q_0&P_0\\Q_1&P_1\end{pmatrix}$, $\det M=-D$. Writing
$(\gamma,\varepsilon)=M^{-1}(a,b)^{T}$ with $(a,b)\in\Z^2$ and using Cramer's rule,
\[
\gamma+\varepsilon\theta=\frac{aP_1-bP_0+\theta(bQ_0-aQ_1)}{-D}=\frac{a\delta_1-b\delta_0}{D},
\]
so $\{\gamma+\varepsilon\theta:(\gamma,\varepsilon)\in\Lambda\}=\LL_D$.

\emph{(A).} Let $z=\gamma+\varepsilon\theta$ with $(\gamma,\varepsilon)\in\Lambda$. Then
$n_j:=\gamma Q_j+\varepsilon P_j\in\Z$ and $Q_jz-n_j=\varepsilon\delta_j$, hence
$a_j\dist{Q_jz}\le|\varepsilon||D|\,a_j/Q_{j+1}\le|\varepsilon D|/Q_j\to0$, using $Q_{j+1}\ge a_jQ_j$.

\emph{(B).} We prove the contrapositive: suppose $\limsup_j a_j\dist{Q_jz}<c<1$; we show $z\in\LL_D$.
Let $n_j$ be the nearest integer to $Q_jz$ and $y_j=Q_jz-n_j$, so $|y_j|<c/a_j$ for $j\ge J_1$.
Put $m_j:=n_{j+1}-a_jn_j-n_{j-1}\in\Z$. Since $Q_j z$ satisfies the recurrence,
$y_{j+1}-a_jy_j-y_{j-1}=-m_j$, so
\[
|m_j|\le|y_{j+1}|+a_j|y_j|+|y_{j-1}|<c\Bigl(1+\tfrac1{a_{j+1}}+\tfrac1{a_{j-1}}\Bigr)<1
\]
for large $j$ (because $a_j\to\infty$ and $c<1$). Hence $m_j=0$ for $j\ge J$. So $(y_j)_{j\ge J-1}$
satisfies the recurrence, and by (e) $y_j=\alpha Q_j+\beta\delta_j$ for $j\ge J-1$ with real
$\alpha,\beta$. Since $y_j\to0$, $\delta_j\to0$ and $Q_j\to\infty$, we get $\alpha=0$. Therefore
$n_j=Q_jz-\beta\delta_j=(z-\beta\theta)Q_j+\beta P_j$ for $j\ge J-1$. These are integers at
consecutive indices, so by Lemma 1 $(z-\beta\theta,\beta)\in\Lambda$, and $z=(z-\beta\theta)+\beta\theta\in\LL_D$.

\emph{(C).} $w_j\dist{Q_jz}=(w_j/a_j)\cdot a_j\dist{Q_jz}$. A $\limsup$ of the second factor that
is $\ge1$ (possibly $\infty$) is multiplied by $w_j/a_j\to\kappa>0$ to give a $\limsup\ge\kappa$,
and a null sequence stays null. \qedhere
\end{proof}

\begin{remark}
Not proved here: the behaviour when $a_j\not\to\infty$; that a single index (rather than the
$\limsup$) suffices; and that the threshold $1$ is optimal for general $(a_j)$.
Its sharpness in one case is Theorem~\ref{thm:sharp}.
\end{remark}

\section{The lower bound $1/4$}\label{sec:lower}

\begin{remark}[Input]\label{rem:input}
We use that $\lim_jP_j/Q_j=e$ for the sequences of the Introduction. This is proved in
Proposition~\ref{prop:input}(i) below; the proof there does not use any result of this section,
so there is no circularity. Independently, the same limit is machine-checked in the accompanying
Lean development (\texttt{Psel\_div\_Qsel\_tendsto\_exp}, via the same Hermite-type integral
representation of $Q_je-P_j$; axioms $[\texttt{propext},\ \texttt{Classical.choice},\ \texttt{Quot.sound}]$
only, as printed in CI run 45759959535; see the verification log).
\end{remark}

\begin{corollary}[Parts (a) and (b) of Theorem~\ref{thm:main}]\label{cor:lower}
Take $a_j=4j+2$, $(Q_0,Q_1,P_0,P_1)=(1,1,1,3)$ and $w_j=j$. Then (H) holds with $D=1\cdot1-3\cdot1=-2$,
$\theta=e$ (Proposition~\ref{prop:input}(i)), $\kappa=\lim j/(4j+2)=1/4$, and $\LL_D=\LL$.
Hence $z\in\LL\Rightarrow j\dist{Q_jz}\to0$ and $z\notin\LL\Rightarrow\limsup_jj\dist{Q_jz}\ge1/4$.
\end{corollary}

\begin{proof}
Here $\delta_0=e-1$, $\delta_1=e-3$, so $\LL_D=\bigl(\Z(e-1)+\Z(e-3)\bigr)/(-2)$. This group contains
$\frac{e-1}2-\frac{e-3}2=1$ and $\frac{e-1}2$, and is generated by $\frac{e-1}2$ and $\frac{e-3}2=\frac{e-1}2-1$,
so $\LL_D=\Z+\Z\frac{e-1}{2}=\Z+\Z\frac{e+1}{2}=\LL$. Apply Proposition~\ref{prop:dich}(A)--(C). \qedhere
\end{proof}

\section{Sharpness: the explicit $z^*$}\label{sec:sharp}

Define integers $n_0=n_1=0$ and for $j\ge1$
\[
n_{j+1}=(4j+2)n_j+n_{j-1}+m_j,\qquad m_j=\begin{cases}1&j\ge2\text{ even},\\0&\text{otherwise,}\end{cases}
\]
and put $z^*:=\lim_{j\to\infty}n_j/Q_j$ (the limit exists by the proof below).

\begin{theorem}[Part (c) of Theorem~\ref{thm:main}]\label{thm:sharp}
$z^*\notin\LL$, $\;j\dist{Q_jz^*}\to1/4$ along even $j$, and $j\dist{Q_jz^*}\to0$ along odd $j$.
In particular $\limsup_jj\dist{Q_jz^*}=1/4$, so $\kappa=1/4$ and the infimum is attained.
\end{theorem}

\begin{proof}
\emph{Growth facts.} (G1) For $j\ge1$, $(4j+2)Q_j\le Q_{j+1}\le(4j+3)Q_j$ (the upper bound since
$Q_{j-1}\le Q_j$). In particular $Q_{j+1}\ge6Q_j$ for $j\ge1$. (G2) $Q_{j+1}\ge(4j+2)(4j-2)Q_{j-1}$ for
$j\ge2$. (G3) For $i_0\ge1$, $\sum_{i\ge i_0}1/Q_{i+1}\le\frac65/Q_{i_0+1}$ (ratio at least $6$).

\emph{Step 1.} Let $E_j:=n_{j+1}Q_j-n_jQ_{j+1}$. Substituting the recurrences,
$E_j=-E_{j-1}+m_jQ_j$, with $E_0=0$. Hence $E_j=(-1)^jS_j$ where
$S_j=\sum_{i\text{ even},\,2\le i\le j}Q_i$ (and $S_0=S_1=0$): indeed $S_j=S_{j-1}+(-1)^jm_jQ_j$.

\emph{Step 2.} For even $i\ge4$, $Q_i/Q_{i-2}\ge6\cdot6=36$ by (G1). Hence if $j'$ is the largest even index
$\le j$ (and $j\ge2$), $Q_{j'}\le S_j\le\frac{36}{35}Q_{j'}$.

\emph{Step 3.} $\dfrac{n_{j+1}}{Q_{j+1}}-\dfrac{n_j}{Q_j}=\dfrac{E_j}{Q_jQ_{j+1}}$ and
$|E_j|/(Q_jQ_{j+1})\le\frac{36}{35}/Q_{j+1}$, which is summable, so $z^*$ exists. With
$t_i:=S_i/(Q_iQ_{i+1})$,
\[
y_j:=Q_jz^*-n_j=Q_j\sum_{i\ge j}(-1)^it_i .
\]

\emph{Step 4 (tail).} By Step 2, $t_i\le\frac{36}{35}/Q_{i+1}$, so by (G3)
$\sum_{i>j}t_i\le\frac{36}{35}\cdot\frac65\cdot\frac1{Q_{j+2}}\le\frac{1.24}{Q_{j+2}}$.

\emph{Step 5 (even $J\ge4$).} Here $t_J=(Q_J+S_{J-2})/(Q_JQ_{J+1})$ and $S_{J-2}\le1.03Q_{J-2}$.
The alternating tail is bounded by $\sum_{i>J}t_i$, so
\[
y_J=\frac{Q_J}{Q_{J+1}}(1+\eta_J),\qquad
|\eta_J|\le\frac{1.03\,Q_{J-2}}{Q_J}+\frac{1.24\,Q_{J+1}}{Q_{J+2}}
\le\frac{1.03}{(4J-2)(4J-6)}+\frac{1.24}{4J+2}\longrightarrow0,
\]
using (G1), (G2). Since $Q_J/Q_{J+1}\in[\frac1{4J+3},\frac1{4J+2}]$, $y_J\in(0,\tfrac12)$ for large even $J$,
so $\dist{Q_Jz^*}=y_J$ and $J\dist{Q_Jz^*}\to1/4$.

\emph{Step 6 (odd $j$).} Now $S_j=S_{j-1}\le1.03Q_{j-1}$, so
\[
|y_j|\le Q_jt_j+Q_j\sum_{i>j}t_i\le\frac{1.03\,Q_{j-1}}{Q_{j+1}}+\frac{1.24\,Q_j}{Q_{j+2}}=O(j^{-2}),
\]
whence $j\dist{Q_jz^*}\le j|y_j|\to0$.

\emph{Step 7.} By Steps 5 and 6 the $\limsup$ is $1/4>0$. If $z^*\in\LL$, Theorem~\ref{thm:main}(a)
would give $j\dist{Q_jz^*}\to0$; so $z^*\notin\LL$. Together with Corollary~\ref{cor:lower}, $\kappa=1/4$ and
it is attained at $z^*$. \qedhere
\end{proof}

\begin{remark}[Numerics, not part of the proof]
The project log records $z^*\approx0.0140416395\allowbreak 6694902840\allowbreak 5640892454$, and
$J\dist{Q_Jz^*}=0.2439$ at $J=20$, $0.24793$ at $J=60$, with odd-index values of
order $j^{-2}$. These were computed by project scripts and were \emph{not} recomputed for this draft;
they are finite computations and prove nothing about the limit.
\end{remark}

\section{The limit $P_j/Q_j\to e$ and the continued fraction of $(e-1)/2$}\label{sec:input}

We now prove the input used in Corollary~\ref{cor:lower}, together with the identification of $Q_j$
as continued-fraction denominators. Write $\delta_j:=Q_je-P_j$ as in Proposition~\ref{prop:dich}.

\begin{proposition}\label{prop:input}
Let $Q_j,P_j$ be as in the Introduction, and for $n\ge0$ put
\[
J_n:=\frac1{n!}\int_0^1x^n(1-x)^n e^x\,dx .
\]
\begin{enumerate}
\item[(i)] For every $j\ge0$, $\delta_j=Q_je-P_j=(-1)^jJ_j$ and $0<J_j\le\dfrac{e-1}{4^j\,j!}$. In particular
$\delta_j\ne0$, $\delta_j\to0$ and $P_j/Q_j\to e$.
\item[(ii)] Let $c_0=0$, $c_1=1$ and $c_k=4k-2$ for $k\ge2$. Then
\[
\frac{e-1}{2}=[c_0;c_1,c_2,\dots]=[0;1,6,10,14,\dots]
\]
is the regular continued-fraction expansion of $(e-1)/2$, its convergents $p_j/q_j$ satisfy
$q_j=Q_j$ and $2p_j=P_j-Q_j$ for all $j\ge0$. Equivalently, $\frac2{e-1}=[1;6,10,14,\dots]$.
\end{enumerate}
\end{proposition}

\begin{proof}
\emph{(i) A recurrence for $J_n$.} Put $g(x)=x(1-x)$ and $f_n=g^n/n!$, so $J_n=\int_0^1f_ne^x\,dx$.
We have $g'=1-2x$, $g''=-2$ and $(g')^2=1-4x+4x^2=1-4g$. For $n\ge1$,
$f_{n+1}'=g^ng'/n!$ and
\[
f_{n+1}''=\frac{n\,g^{n-1}(g')^2+g^ng''}{n!}
=\frac{g^{n-1}}{(n-1)!}-\frac{4n\,g^n}{n!}-\frac{2g^n}{n!}=f_{n-1}-(4n+2)f_n .
\]
Since $n+1\ge2$, both $f_{n+1}$ and $f_{n+1}'$ vanish at $x=0$ and $x=1$ (each contains the factor $g^n$
with $n\ge1$). Integrating by parts twice,
\[
\int_0^1f_{n+1}''e^x\,dx=\bigl[f_{n+1}'e^x\bigr]_0^1-\bigl[f_{n+1}e^x\bigr]_0^1+\int_0^1f_{n+1}e^x\,dx
=\int_0^1f_{n+1}e^x\,dx .
\]
Hence $J_{n+1}=J_{n-1}-(4n+2)J_n$ for $n\ge1$. Directly, $J_0=\int_0^1e^x\,dx=e-1$ and, using
$\int_0^1xe^x\,dx=1$ and $\int_0^1x^2e^x\,dx=e-2$, $J_1=1-(e-2)=3-e$.

\emph{Identification with $\delta_j$.} Put $u_j:=(-1)^jJ_j$. Then $u_0=e-1=\delta_0$, $u_1=e-3=\delta_1$, and for
$n\ge1$, $u_{n+1}=(-1)^{n+1}\bigl(J_{n-1}-(4n+2)J_n\bigr)=u_{n-1}+(4n+2)u_n$. The sequence
$\delta_j=Q_je-P_j$ satisfies the same recurrence (it is a fixed linear combination of $(Q_j)$ and $(P_j)$),
so $\delta_j=u_j$ for all $j\ge0$ by induction.

\emph{Bounds.} On $(0,1)$ the integrand of $J_j$ is positive and $0<g\le\frac14$, so
$0<J_j\le\frac1{4^jj!}\int_0^1e^x\,dx=\frac{e-1}{4^jj!}$. Since $Q_j\ge1$,
$\bigl|P_j/Q_j-e\bigr|=|\delta_j|/Q_j\le\frac{e-1}{4^jj!}\to0$.

\emph{(ii) The continued fraction.} Let $p_{-1}=1$, $p_0=c_0=0$, $q_{-1}=0$, $q_0=1$ and, for $k\ge1$,
$p_k=c_kp_{k-1}+p_{k-2}$, $q_k=c_kq_{k-1}+q_{k-2}$; these are the convergents $p_k/q_k=[c_0;c_1,\dots,c_k]$.
Then $q_0=1=Q_0$, $q_1=c_1q_0+q_{-1}=1=Q_1$, and for $k\ge1$ we have $c_{k+1}=4k+2$, so
$q_{k+1}=(4k+2)q_k+q_{k-1}$: this is the recurrence of $(Q_j)$, hence $q_k=Q_k$ for all $k\ge0$. Similarly
$p_0=0$, $p_1=1$, and $r_k:=Q_k+2p_k$ satisfies $r_{k+1}=(4k+2)r_k+r_{k-1}$ for $k\ge1$ with
$r_0=1=P_0$ and $r_1=3=P_1$; hence $P_k=Q_k+2p_k$ for all $k\ge0$. Therefore, by (i),
\[
\frac{p_k}{q_k}=\frac12\Bigl(\frac{P_k}{Q_k}-1\Bigr)\longrightarrow\frac{e-1}2 .
\]
Thus the infinite regular continued fraction $[c_0;c_1,c_2,\dots]$ (whose value is by definition
$\lim_kp_k/q_k$) equals $(e-1)/2$. The number $(e-1)/2$ is irrational (by Proposition~\ref{prop:dich}(d)
with $\theta=e$, or directly from (i)), and an irrational number has a unique infinite regular
continued-fraction expansion: if $x=[c_0;c_1,\dots]$ with $c_k\ge1$ for $k\ge1$, then the complete quotients
$x_k=[c_k;c_{k+1},\dots]$ satisfy $x_k=c_k+1/x_{k+1}$ and $x_{k+1}>c_{k+1}\ge1$, so $c_k=\lfloor x_k\rfloor$
is forced by the Gauss algorithm. Hence $[0;1,6,10,14,\dots]$ is \emph{the} regular continued-fraction
expansion of $(e-1)/2$, with denominators $q_j=Q_j$. Finally, $\frac2{e-1}=1/\frac{e-1}2=[c_1;c_2,c_3,\dots]=[1;6,10,14,\dots]$. \qedhere
\end{proof}

\begin{remark}[Finite checks, not part of the proof]
The polynomial identity $f_{n+1}''=f_{n-1}-(4n+2)f_n$ (exact rational arithmetic, $1\le n\le40$),
the relations $q_j=Q_j$ and $P_j=Q_j+2p_j$ ($0\le j<400$), the identity $\delta_j=(-1)^jJ_j$ (numerical
quadrature at $300$ digits, $0\le j\le30$, maximal relative error $2.5\times10^{-201}$), and the first $80$
partial quotients of $(e-1)/2$ computed by the Gauss algorithm at $3000$ digits all agree with
Proposition~\ref{prop:input}. Script and output: \texttt{logs/verification/R5/input\_prop\_chk.py}
and \texttt{.out}. These are finite computations.
\end{remark}

\section{Formal verification}\label{sec:lean}

The core of the argument has been formalized in Lean~4. The toolchain is
\texttt{leanprover/lean4:v4.35.0}, with Mathlib at the matching release \texttt{v4.35.0}.
We list exactly what has been checked.

\medskip\noindent\textbf{Sharpness.} File \texttt{EpiFormal/KappaSharp.lean}. The theorem
\texttt{kappa\_eq\_quarter} states: (i) for every $z\notin\LL$ and every $c<1/4$,
it is \emph{not} the case that $j\dist{Q_jz}\le c$ eventually; and (ii) there exists $z\notin\LL$ such that
for every $c>1/4$, $j\dist{Q_jz}\le c$ eventually. Supporting results are
\texttt{KappaSharp.upper\_ev}, \texttt{KappaSharp.lower\_freq} and \texttt{kappa\_sharp}; the last
yields the witness (built from the same defect sequence as $z^*$) with $j\dist{Q_jz}\ge1/5$ frequently.
Together they give $\kappa=1/4$, attained. Evidence: GitHub Actions workflow
\texttt{lean-audit.yml}, run id \texttt{45759959535}, head commit
\texttt{1e9b94cf6fb8a6df5f618505f8888d45e6a611ff} (PR \#4), conclusion success; full \texttt{lake build}
(``Build completed successfully (9392 jobs)''); no \texttt{sorry}/\texttt{admit} in the file. A native local
compile was not completed, so CI is the build of record. \texttt{\#print axioms} for
\texttt{upper\_ev}, \texttt{lower\_freq}, \texttt{kappa\_sharp}, \texttt{kappa\_eq\_quarter}:
\[
[\texttt{propext},\ \texttt{Classical.choice},\ \texttt{Quot.sound}]
\]
i.e.\ Lean foundational axioms (\texttt{propext}, \texttt{Quot.sound}) plus the classical-whitelist
axiom \texttt{Classical.choice}; no project-added axioms. A statement-fidelity review of
\texttt{kappa\_eq\_quarter} against the informal statement was recorded as passed (an internal review,
not an external one).

\medskip\noindent\textbf{Lower bound and part (a).} The declarations
\texttt{Classification.corollary3\_gap} (the lower bound for $z\notin\LL$, with every $c<1/4$),
\texttt{Classification.theorem2\_classification}
($j\dist{Q_jx}\to0\iff x\in\LL$) and \texttt{Classification.theorem2\_sufficiency} were re-run
natively (\texttt{lake env lean} on the audit file, exit code $0$) with the same axiom list
$[\texttt{propext},\ \texttt{Classical.choice},\ \texttt{Quot.sound}]$ and no project axioms.
All three were rebuilt in the full CI \texttt{lake build} (run 45759959535, 9392 jobs, success), and their axiom lists printed there are the same. The CI exact-match axiom gate covers \texttt{corollary3\_gap} and \texttt{theorem2\_classification}; for \texttt{theorem2\_sufficiency} the list is printed but not gated.

\medskip\noindent\textbf{Not machine-checked.}
\begin{itemize}
\item The finer statement that $j\dist{Q_jz^*}\to1/4$ along even $j$ and $\to0$ along odd $j$
(Theorem~\ref{thm:sharp}) is established on paper only; Lean covers the $\limsup/\liminf$-type
consequences listed above.
\item The general Proposition~\ref{prop:dich} (arbitrary $a_j$) is not formalized; Lean treats the
sequence $a_j=4j+2$ only.
\item The identification of $Q_j$ with the continued-fraction denominators of $(e-1)/2$
(Proposition~\ref{prop:input}(ii)) is proved on paper only; in Lean, $Q_j$ is a recurrence definition
and the identification is not a theorem there.
\item The limit $P_j/Q_j\to e$ is proved on paper in Proposition~\ref{prop:input}(i) and, independently,
machine-checked in Lean (\texttt{Psel\_div\_Qsel\_tendsto\_exp}, Remark~\ref{rem:input}). The paper proof
of Proposition~\ref{prop:input} itself has not been compared line by line with the Lean proof.
\end{itemize}
The Lean definitions \texttt{Qsel} ($Q_0=Q_1=1$, $Q_{j+1}=(4j+2)Q_j+Q_{j-1}$, written with a
shifted index), \texttt{distToNearestInt} ($x\mapsto|x-\operatorname{round}x|$) and \texttt{latticeL}
($\mathbb Z+\mathbb Z(e+1)/2$) were checked against this note by an independent internal review,
which also confirmed that the Lean witness equals $z^*$. An external reader should re-check this
before submission (TODO before submission).

\section{Scope and limitations}\label{sec:scope}

\begin{itemize}
\item \textbf{Nothing about $\pi$ or $e+\pi$.} Theorem~\ref{thm:main} holds for every real $z$ and gives no
information about any particular number unless its membership in $\LL$ is known. No statement about
the irrationality or transcendence of $e+\pi$ (or $\pi$, $e\pi$, etc.) follows or is claimed.
\item \textbf{Sharpness is for one sequence.} Optimality of $1/4$ is proved for $a_j=4j+2$ and
$(Q_0,Q_1)=(1,1)$ only. The threshold $1$ in Proposition~\ref{prop:dich}(B) is not shown to be sharp in general.
\item \textbf{Qualitative content is not new.} By the identification with continued-fraction
denominators of $(e-1)/2$ (Proposition~\ref{prop:input}(ii)), parts (a) and (b) in qualitative form follow
from \cite{KL}. The expansion $(e-1)/2=[0;1,6,10,\dots]$ itself is not claimed as new; we have not checked
a classical source for it and prove it only to keep the note self-contained. At most the explicit
constant $1/4$, the lower bound for general weights in Proposition~\ref{prop:dich}(C), and the attainment
example are candidates for new content.
\item \textbf{Novelty rests on a negative search.} We found no precedent for the value
$\kappa=1/4$ with a sharpness example, but this is a \emph{negative search}, not a proof that none exists.
Coverage: full texts of \cite{KL}, \cite{Larcher} and the open-access copy of \cite{BDS}; KL's paraphrase
of \cite{Stewart}; a further metric paper of Philipp (not used here). Not reached: full text of
\cite{Stewart} (paywalled or access-blocked); Eggleston and Erd\H{o}s--Taylor (not opened);
Cusick--Flahive and Rockett--Sz\"usz (not reached). Hence an earlier appearance in these sources
cannot be excluded.
\item \textbf{Numerics.} The numerical values in Sections~\ref{sec:sharp} and~\ref{sec:input} and in the
comparison remark are finite computations; they are not used in any proof.
\item \textbf{Formalization.} See Section~\ref{sec:lean} for the exact boundary of what is machine-checked.
\end{itemize}

% TODO before submission: human novelty review (obtain Stewart full text), author check of all proofs,
% (P_j/Q_j -> e input: now proved in prop:input, R5), optionally add a full-text-checked classical
% reference for (e-1)/2 = [0;1,6,10,...], decide on journal, add acknowledgements/funding statement.

\begin{thebibliography}{9}

\bibitem{KL} Kraaikamp, Liardet, \emph{Good approximations and continued fractions},
Proc.\ Amer.\ Math.\ Soc.\ \textbf{112} (1991), no.~2, 303--309.
DOI 10.1090/S0002-9939-1991-1062392-7 (DOI from Crossref; not printed on the PDF read).

\bibitem{Larcher} Larcher, \emph{A convergence problem connected with continued fractions},
Proc.\ Amer.\ Math.\ Soc.\ \textbf{103} (1988), no.~3, 718--722.

\bibitem{BDS} B\'ir\'o, Deshouillers, S\'os,
Studia Sci.\ Math.\ Hungar.\ \textbf{38} (2001), 97--113.
DOI 10.1556/sscmath.38.2001.1-4.7. (Open-access copy read at \texttt{http://real.mtak.hu/110616/}.)

\bibitem{Stewart} M. Stewart, \emph{Irregularities of uniform distribution},
Acta Math.\ Acad.\ Sci.\ Hungar.\ \textbf{37} (1981), 185--221.
DOI 10.1007/bf01904883 (Crossref). Cited only via the reference list and Remark~3 of \cite{KL};
\emph{full text not consulted}. (Another source cites different page numbers; unresolved.)

\end{thebibliography}

\end{document}
